\section{Why the MIP does not scale}
\label{sec:mip}

We give five reasons. They add up: the first two concern the sequential solve, the other three
concern parallel and federated solves.

\paragraph{1. Size of the formulation.}
The \exdbn{} MIQP has $d(d-1)+pd^2$ binary indicators, the same number of continuous
weights, twice as many big-$M$ constraints, and $\binom d2$ two-cycle constraints. The family
of cycle constraints is exponential: the complete digraph on $d$ nodes has
$\sum_{k=2}^{d}\binom dk (k-1)!$ simple cycles, so they can only be generated lazily, and each
incumbent may trigger a separation round. The row-wise objective used by the reference
implementation has $O(nd\,m^2)$ quadratic terms before aggregation. For $n=1000$ and $p=1$ that is
about $2.5\cdot 10^8$ terms for $d=50$ and $2\cdot 10^9$ for $d=100$. Building the $d=50$ model took
18 minutes and 12.3\,GB of memory on our test machine (\cref{sec:exp-scaling}). The $d=100$
model has eight times as many terms and would exceed its 32\,GB. The Gram form of \cref{prop:gram} removes the factor $n$ but leaves the
combinatorial core untouched.

\paragraph{2. Worst-case and typical complexity.}
Score-based structure learning is NP-hard \citep{chickering2004large}, so branch-and-bound
must be expected to explore exponentially many nodes in the worst case. The typical case is
also hard here: the big-$M$ relaxation is very weak. With $M=100$ and data-scale weights of
order one, the relaxation can set $e_{ij}=|w_{ij}|/M\approx 10^{-2}$. This pays a negligible
fraction of the penalty $\lambda/d$, and it satisfies every cycle inequality
$\sum_{C}e_{ij}\le|C|-1$ trivially. The root bound is therefore essentially the
\emph{unconstrained} least-squares loss with no acyclicity and no sparsity, and the gap can
only be closed by branching. Our experiments show exactly this: the MIP gap at the time limit
grows quickly with $d$ (\cref{sec:experiments}).

\paragraph{3. Even the root relaxation is inherently sequential.}
Suppose we gave up exactness and solved only the root node, the continuous relaxation, a
convex QP that contains linear programming as a special case. Linear programming is
P-complete under log-space reductions \citep{dobkin1979lp}, and remains P-complete even when
only an approximate solution is sought \citep{serna1991approximating}. Unless
$\mathrm{P}=\mathrm{NC}$, that is unless every polynomial-time problem can be solved in
polylogarithmic time on polynomially many processors (Nick's class), there is no algorithm
that solves general LPs in polylogarithmic depth \citep{greenlaw1995limits}. Our
relaxations are a structured subclass, and we do not claim that this particular family is
P-complete. But no structure that would place it in NC is known, and general-purpose
solvers treat it as a general LP or QP. In practice, interior-point methods need
$O(\sqrt{N}\log(1/\varepsilon))$ \emph{sequential} Newton steps, each a sparse factorisation
that parallelises only moderately. First-order LP methods \citep{applegate2021pdlp}
parallelise each iteration well, but need a number of iterations that grows with condition
measures of the instance. By contrast, one iteration of \feddbn{} consists of matrix products,
which have logarithmic depth, plus averaging, and the number of rounds is a parameter chosen
in advance. The greedy order projection is sequential in $d$, but it does $O(d^2)$ work, uses
no data, and runs once per epoch.

\paragraph{4. Parallel branch-and-bound scales poorly in practice.}
Branch-and-bound looks embarrassingly parallel, yet the empirical record is sobering.
ParaSCIP \citep{shinano2010parascip} solved two previously open MIPLIB2003 instances with up
to 2048 cores of the HLRN~II supercomputer. This took about 86 and 114 hours, with 16 and 10
checkpoint restarts. The ramp-up phase alone grew from roughly 250\,s on 512 cores to 535\,s
on 2048 cores. The authors note that most earlier parallelisation approaches ``struggled to
catch up with the performance of state-of-the-art'' sequential solvers, and that many open
instances were first solved sequentially. They also report that longer ramp-up and ramp-down
phases with more cores hamper scalability. \citet{koch2012million} discuss the
obstacles to using $10^6$ cores on a single integer program. These include the ramp-up until
the tree is large enough to feed all workers, communication and load balancing, and the fact
that the search, and hence the total work, changes with the degree of parallelism. For a
broader survey of how parallel high-performance computing and emerging architectures are
used for mathematical optimisation, including branch-and-bound, see
\citet{bernalneira2024survey}.

\paragraph{5. Federated branch-and-bound is either centralised or chatty.}
Every node relaxation depends on all the data. There are two options. Either the server
receives the pooled sufficient statistic $S$ (\cref{prop:gram}) and runs a centralised
branch-and-bound, which is subject to reasons 1 to 4. Or every node relaxation is solved
distributedly, say by ADMM across clients, which costs tens to hundreds of communication
rounds \emph{per node}, over millions of nodes. \feddbn{} uses a few hundred rounds in
total.
